\chapter{Аналитическая часть}

В аналитической части будет проведена формализация задачи коммивояжёра, описан классический муравьиный алгоритм и рассмотрены его основные модификации. Будет выполнен сравнительный анализ существующих подходов к решению задачи, выделены их достоинства и недостатки, а также обоснован выбор направления для разработки комбинированного метода.

\section{Формализация задачи коммивояжёра}

Словесно задача коммивояжёра формулируется следующим образом: коммивояжёр, выходящий из некоторого города желает посетить $n - 1$ других городов ровно по 1 разу и вернуться к исходному. Известно, что из каждого города можно перейти в любой другой и также известны расстояния между ними. Требуется установить в каком порядке коммивояжёр должен посещать города, чтобы пройденное им расстояние было минимальным~\cite{mudrov}.

Данное определение возможно формализовать математически, введя несколько определений из теории графов~\cite{belousev}:

\textbf{Неориентированным графом (далее неорграф) $G_n$} называется пара $(V_n, E_n)$, где $V_n$ -- конечное множество мощности $n$, элементы которого называются \textbf{вершинами}, а $E_n$ -- множество неупорядоченных пар $e_{ij} = \{v_i, v_j\}$ на $V_n$, элементы которого называют \textbf{рёбрами}
           
\textbf{Неорграф $G_n$} называется \textbf{полным}, если для любых пар вершин $v_i$, $v_j$ ($i < j$) из множества $V_n$ существует ребро $e_{ij}$ в множестве $E_n$.

\textbf{Конечной цепью} $X_k$ в неорграфе $G_n$ называется произвольная упорядоченная последовательность вершин $v_{i_1}, v_{i_2}, \dots, v_{i_k}$ из множества $V_n$, таких, что $\forall j, 1 <= j < k, e_{i_j, i_{j+1}} = \{v_{i_j}, v_{i_{j + 1}}\} \in E_n$ , т.~е. такие, что соседние вершины пути связаны ребром.

\textbf{Циклом} $Y$ в неорграфе $G_n$ называют такую цепь, у которой первая и последняя вершина совпадают, при этом цикл называют \textbf{гамильтоновым}, если вершины, кроме первой и последней, не повторяются. В дальнейшем гамильтонов цикл будет записываться как  $(v_{i_1}, v_{i_2}, \dots, v_{i_n})$, где $v_{i_1}$ -- первая и последняя вершина цепи. 

\textbf{Взвешенным неорграфом (размеченным)} называется пара $(G_n, F)$, где $F$ -- функция сопоставляющая каждому ребру графа некоторое действительное положительное число, то есть $F: E_n \rightarrow \mathbb{R}$, называемая функцией разметки. При этом такое число называется \textbf{весом} или \textbf{длиной} ребра. Весом цепи $X_k$ назовём сумму весов всех рёбер, входящих в него: $F(X_m) = \sum^{k-1}_{i=1}F(e_{j_i j_{i+1}})$, а весом цикла $Y$ -- $F(X_m) = F(e_{j_m j_1}) + \sum^{k-1}_{i=1}F(e_{j_i j_{i+1}})$.

Тогда, используя вышеизложенную терминологию, задачу коммивояжёра возможно определить следующим образом: Пусть дан полный взвешенный неорграф $G_n$ и  функция разметки над этим графом $F$. Требуется найти гамильтонов цикл с наименьшим (наибольшим) весом (оптимальный), среди всех остальных гамильтоновых циклов данного графа.

На рисунке~\ref{tsp-idef0} представлена формализация задачи коммивояжёра в виде IDEF0-диаграммы.

\begin{figure}[H]
	\centering
	\includegraphics[width=0.9\textwidth]{tsp/01\_A0}
	\caption{Формализация задачи коммивояжёра в виде IDEF0-диаграммы}
	\label{tsp-idef0}
\end{figure}

\section{Классический муравьиный алгоритм}

Классический муравьиный алгоритм впервые был предложен в~\cite{ant-origin} как новый подход к решению задач комбинаторной оптимизации, в частности, как решение задачи коммивояжёра. Метод построен на аналогии с тем, как функционируют муравьи при поиске маршрутов от колонии до источников пищи. Муравьи используют химические маркеры, феромоны, оставляемые на пройденном ими пути, как средство коммуникации. По началу муравьи ходят случайно, тем самым исследуя различные пути к источнику пищи. Пройдя путь в одну сторону и обратно муравьи оставляют феромоны на пройденном пути, причём концентрация оставленных ими на пути обратно зависит от длины или сложности пути. Последующие группы муравьёв смогут улавливать оставленный на путях феромон и с большей вероятностью будут выбирать те пути, на которых его больше, тем самым выбирая более простые и короткие пути.

Муравьиный алгоритм симулирует это поведение муравьёв на графе и итеративно ищет кратчайшие пути. На нулевой итерации происходит инициализация всей системы, а на последующих происходит поиск кратчайшего маршрута через все узлы, через имитацию муравьём с помощью упрощённых агентов.

\subsection{Этап инициализации}

Пусть дан взвешенный неограф ($G_n$, $F$). Тогда $\tau_{ij}(k)$ -- интенсивность феромона на ребре $(v_i, v_j)$ графа $G_n$ на k-ой итерации метода, где $1\leqslant i,j \leqslant n$. Так как граф неориентированный, то $\tau_{ij}(k)~=~\tau_{ji}(k), \forall k$.

На нулевой итерации все ребра не обладают феромоном ($\tau_{ij}(0) = 0$), либо, как предложено в~\cite{ant-origin} обладают очень малым константным значением распылённого феромона ($\tau_{ij}(0) = c, c\geqslant0$).

\subsection{Основной этап}

Основной этап описывает одну итерацию поиска пути оптимального пути. На каждой итерации запускается популяция искусственных агентов -- \textbf{муравьёв}, каждый из которых независимо строит допустимый маршрут. Поведение агента определяется 3-мя его основными свойствами:

\begin{itemize}
	\item \textbf{зрение} -- агент способен оценить расстояние (вес) от одной вершины графа, до
	другой;
	\item \textbf{обоняние} -- агент способен оценить концентрацию феромона, распылённого на каждом ребре графа, исходящим из вершины, где он находиться;
	\item \textbf{память} -- агент способен запоминать те вершины, на которых он уже был и, соответственно в дальнейшем посещать только нет посещённые до этого вершины.
\end{itemize}

\textbf{Построение маршрутов}

Начиная с заданной или случайной вершины графа отдельный муравей итеративно выбирает по одному ребру, по которому перейти. Выбор ребра осуществляется с помощью стохастического правила, которое учитывает посещённые до этого муравьём вершины, распылённый на рёбрах феромон и их вес Конкретное ребро для перехода определяется случайно, с вероятностями, описанными ниже.

Пусть на итерации $k$ муравей $a$ находится в вершине $v_i$, а множество ещё не посещённых им вершин обозначается как $U_a(i)$. Тогда вероятность выбора вершины $v_j \in U_a(i)$ в качестве следующей задаётся формулой:

\begin{equation}
P_{ij}^{(a)}(k) =
\begin{cases}
	\displaystyle \frac{\tau_{ij}(k)^\alpha \cdot \eta_{ij}^\beta}
	{\sum\limits_{v_l \in U_a(i)} \tau_{il}(k)^\alpha \cdot \eta_{il}^\beta}, & \text{если } v_j \in U_a(i); \\
	0, & \text{иначе}, \\
\end{cases}
\end{equation}

где
\begin{itemize}
	\item $\tau_{ij}(k)$ -- интенсивность феромона на ребре $e_{ij}$ на итерации $k$;
	\item $\eta_{ij}$ -- эвристическая значимость ребра;
	\item $\alpha \geqslant 0$ -- параметр, определяющий значимость феромона при выборе пути;
	\item $\beta \geqslant 0$ -- параметр, определяющий значимость эвристики при выборе пути.
\end{itemize}

Эвристическая значимость ребра $\eta_{ij}$ описывает то, насколько ребро привлекательно для муравья без учёта распылённого на нём феромона. Эта оценка может быть задана любой функцией, главное её свойство -- неизменность в течении всей симуляции. В~\cite{ant-origin} это функция определяется как обратно пропорциональная весу ребра: $\eta_{ij} = \dfrac{1}{F(e_{ij})}$.

Параметры $\alpha$ и $\beta$ определяют относительную значимость феромонов и эвристики. При $\alpha = 0$ алгоритм становится полностью зависимым от эвристической функции, т.~е. близким к жадному алгоритму. При $\beta = 0$ алгоритм не учитывает вес ребра, а значит становится полностью случайным (т.~к. на начальной итерации концентрация феромона на всех рёбрах одинакова).



Каждый муравей завершает построение своего тура, когда множество $U_a(i)$ становится пустым, после чего он возвращается в исходную вершину, таким образом формируя гамильтонов цикл.

После того как все муравьи завершили построение своих маршрутов, происходит вычисление длины каждого цикла:

\begin{equation}
L_a(k) = F(Y^k_a),
\end{equation}

где $Y^k_a$ -- цикл, построенный муравьём $a$.

Минимальный среди всех $L_a(k)$ путь на итерации $k$ называют \textbf{локально лучшим}.

\textbf{Обновление феромонов}

Заключительным этапом каждой итерации является обновление концентрации феромонов на каждом из рёбер. Этот процесс включается в себя два этапа.

\textbf{Испарение феромона}

Данный механизм действует как отрицательная обратная связь, призванный сделать менее удачные пути менее значимыми и предотвратить быстрого попадания алгоритма в локальный минимум. Новый уровень феромона для следующей итерации рассчитывается как:

\begin{equation}
	\tau_{ij}(k + 1) = (1 - p)\tau_{ij}(k), 
\end{equation}
где $p \in [0, 1]$ -- коэффициент испарения феромона.

\textbf{Добавление нового феромона}

Муравьи, построившие маршруты, добавляют феромоны на рёбра, по которым он прошли

\begin{equation}
	\tau_{ij}(k+1) = \tau_{ij}(k+1) + \sum_{a=1}^{m} \Delta \tau_{ij}^{(a)}(k), 
\end{equation}

\begin{equation}
	\Delta \tau_{ij}^{(a)}(k) = 
	\begin{cases}
		\displaystyle \frac{Q}{L_a(k)}, & \text{если муравей } a \text{ прошёл через ребро } (i,j), \\[1ex]
		0, & \text{иначе},
	\end{cases}
\end{equation}
где $Q$ -- положительная константа, определяющая масштаб увеличения концентрации феромона.

Поскольку увеличение концентрации обратно зависит от длины пути, то ребра, входящие в более короткие маршруты, получат больше феромона.

\subsection{Критерии остановки}

Итерационный процесс продолжается до выполнения одного из стандартных критериев:

\begin{itemize}
	\item достижение заранее заданного числа итераций;
	\item отсутствие улучшений локального или глобального лучшего маршрута в течение некоторого количества итераций;
	\item достижение заданного порога качества решения.
\end{itemize}

На выходе алгоритма фиксируется лучший найденный гамильтонов цикл, который и рассматривается как приближённое решение задачи коммивояжёра.

\subsection{Формализация муравьиного алгоритма}

На рисунках~\ref{aco-a0}-\ref{aco-a1} представлена формализация муравьиного алгоритма в виде диаграмм в нотации IDEF0 1-го и 2-го уровня

\begin{figure}[H]
	\centering
	\includegraphics[width=0.9\textwidth]{aco/01\_A0}
	\caption{Формализация муравьиного алгоритма в виде IDEF0-диаграммы 1-го уровня}
	\label{aco-a0}
\end{figure}

\begin{figure}[H]
	\centering
	\includegraphics[width=0.9\textwidth]{aco/02\_A0}
	\caption{Формализация муравьиного алгоритма в виде IDEF0-диаграммы 2-го уровня}
	\label{aco-a1}
\end{figure}


\section{Модификации муравьиного алгоритма}

Как уже было указано во введении, классический муравьиный алгоритм обладает рядом существенных недостатков, обусловленных его структурой~\cite{impaco}:

\begin{itemize}
	\item \textbf{Медленная сходимость.} На начальных итерациях феромон либо отсутствует, либо распределён равномерно, поэтому перемещение муравьёв почти случайно. Различия в концентрации феромона накапливаются медленно из-за испарения и случайности выбора переходов, что приводит к медленному формированию предпочтительных маршрутов.
	
	\item \textbf{Попадание в локальные минимумы.} Если несколько муравьёв случайно усилят подоптимальный маршрут, концентрация феромона на его рёбрах быстро возрастает, и остальные агенты начинают повторять этот путь. Из-за формулы вероятностного выбора возникает эффект самоусиления, приводящий к преждевременной фиксации не оптимального решения.
	
	\item \textbf{Склонность к застою.} При доминировании влияния феромона алгоритм концентрируется на одном маршруте, альтернативы быстро теряют феромон и перестают рассматриваться. В результате муравьи фактически перестают исследовать пространство решений, и процесс оптимизации останавливается.
	
	\item \textbf{Сильная зависимость от параметров.} Эффективность работы алгоритма существенно зависит от выбора параметров $\alpha$, $\beta$, $p$ и $Q$. Небольшие изменения их значений значительно смещают баланс между исследованием и использованием накопленной информации, что влияет на стабильность и качество результата.
	
	\item \textbf{Высокая вычислительная сложность.} На каждой итерации каждый муравей строит полный маршрут, что требует порядка $O(n^2)$ операций, а также производится обновление феромона на всех рёбрах графа $O(n^2)$. В совокупности это делает одну итерацию алгоритма достаточно трудоёмкой.
\end{itemize}

Модификации муравьиного алгоритма направлены на устранение описанных выше проблем и достигаются за счёт изменения одной или нескольких структурных частей классического метода:

\begin{itemize}
	\item \textbf{стратегии выбора пути} -- изменение правил выбора следующей вершины, использование локальных пост- и пред-оптимизаций, дополнительных ограничений;
	
	\item \textbf{правил обновления феромона} -- модификация схем испарения и обновления феромона;
	
	\item \textbf{критериев остановки} -- изменение условий остановки на учитывающие процесс поиска решения.
\end{itemize}


В рамках данной работы выделены три большие группы существующих модификаций муравьиного алгоритма:

\begin{enumerate}
	\item \textbf{расширения классического муравьиного алгоритма} -- модификации, изменяющие отдельные структурные компоненты базового метода;
	
	\item \textbf{гибридные алгоритмы} -- подходы, сочетающие муравьиный алгоритм с другими эвристическими или метаэвристическими методами;
	
	\item \textbf{алгоритмы с динамической параметризацией} -- методы, в которых параметры $\alpha$, $\beta$, $p$, $Q$ и другие регулируются динамически в процессе работы.
\end{enumerate}

Далее будет подробно рассмотрена каждая из групп.

\subsection{Расширения муравьиного алгоритма}

Расширения представляют собой модификации классического муравьиного алгоритма, которые сохраняют его базовую структуру, но изменяют отдельные стратегии, например выбор маршрута, обновления феромона или добавляют дополнительные управляющие конструкции.

\subsection*{Элитарный муравьиный алгоритм (ЭМА)}

Идея элитарных муравьёв основана на концепте элитарности из других эволюционных алгоритмов, например генетическом алгоритме, где лучшие особи (представители) поколения или всех поколений сохраняются на протяжении длительного числа итераций или влияют на процесс поиска решений~\cite{dorigo-book}. Ключевая идея метода в том, чтобы усилить влияние лучшего найденного маршрута на формирование феромонов, тем самым увеличивая вероятность того, что муравью будут включать участки лучшего пути в свои маршруты, что должно повышать сходимость алгоритма.

Модификация изменяет процесс изменения феромонов, добавляя элитных муравьёв, которые всегда ходят по лучшему найденному маршруту в текущей итерации или всех итерациях. Таким образом общая формула изменения феромонов становится:

\begin{equation}
	\tau_{ij}(k+1)=(1 - \rho)\tau_{ij}(k) + \sum_{a=1}^{m} \Delta\tau_{ij}^{(a)} + e \cdot \Delta\tau_{ij}^{best},
\end{equation}


где
\begin{itemize}
	\item $e$ -- количество элитных муравьёв;
	\item $\Delta\tau_{ij}^{best} = \frac{Q}{L_{\text{best}}}$, если ребро $(i,j)$ принадлежит лучшему пути, иначе $0$;
	\item $L_{best}$ -- вес лучшего найденного решения в поколении или за все поколения;
\end{itemize}

\textbf{Улучшения относительно классического алгоритма}

\begin{itemize}
	\item ускоренная сходимость за счёт усиленной эксплуатации лучшего пути;
	\item уменьшение количества случайных отклонений после того, как найден хороший маршрут.
\end{itemize}

\textbf{Ухудшения относительно классического алгоритма}

\begin{itemize}
	\item увеличенный риск падения в локальный минимум и стагнации.
\end{itemize} 

\subsection*{Макс-мин муравьиный алгоритм (МММА)}

Макс-мин муравьиный алгоритм, предложенный в~\cite{mmas} стремится улучшить два момента в классическом муравьином алгоритме:

\begin{enumerate}
	\item увеличить сходимость, за счёт эффективного использования лучших найденных решений;
	\item уменьшить вероятность падения в локальный минимум.
\end{enumerate}

Эти улучшения достигаются за счёт ряда изменений, касающихся управления феромоном:
\begin{enumerate}
	\item в каждой итерации только один муравей распыляет феромон на его пути. Это может быть как муравей, с лучшим путём в текущей итерации, так и муравей, с лучшим путём за все итерации: $$\tau_{ij}(k+1)=(1 - \rho)\tau_{ij}(k) + \Delta\tau_{ij}^{best}$$ Распыление феромона только на одном, лучшем пути выделяет его на фоне остальных, что, как и в элитарном варианте, увеличивает шанс включения его участков в пути других муравьёв;
	\item возможные значения концентрации феромона на всех рёбрах ограничено отрезком $[\tau_{min}, \tau_{max}]$, задаваемых как параметры: $$\tau_{ij}(k+1) = min(\tau_{max}, max(\tau_{min}, \tau_{ij}(k+1)))$$ Введение ограничений на уровни феромона позволяют избежать стагнации и падения в локальный минимум, так как с одной стороны все рёбра будут иметь хоть какое-то количества феромона, что даёт шанс их выбора, с другой не даёт локально лучшим путям получить сверх высокие значения феромона;
	\item начальное значение феромона на рёбрах устанавливается на $\tau_{max}$, что увеличивает вероятность исследования в начале алгоритма.
\end{enumerate}

\textbf{Улучшения относительно классического алгоритма}

\begin{itemize}
	\item уменьшение риска падения в локальный минимум, за счёт ограничения феромона;
	\item ускорение сходимости, за счёт использования лучших путей;
	\item уменьшение операций распыления феромона.
\end{itemize}

\textbf{Ухудшения относительно классического алгоритма}

\begin{itemize}
	\item увеличенная сложность параметризации, чувствительность к выбору границ.
\end{itemize} 

\subsection*{Ранговый муравьиный алгоритм (РМА)}

Ранговый муравьиный алгоритм~\cite{rbas} представляет собой дальнейшее развитие элитарного муравьиного алгоритма, в котором ещё больше взято от эволюционных алгоритмов, в частности генетического алгоритма. В отличие от классического муравьиного алгоритма, в данном подходе при распылении феромона участвуют не все муравьи, а только $w$ лучших.

Муравьи упорядочиваются по возрастанию длины построенного пути:

$$L_1(k)\leqslant L_2(k)\leqslant \dots\leqslant L_m(k),$$

где $L_1(k)$ -- лучший маршрут.

Каждый из первых $w$ лучших муравьёв вносит вклад в феромон пропорционально своему рангу:

\begin{equation}
	\Delta \tau_{ij}^{(r)}(k) = 
	\begin{cases}
		\displaystyle \frac{w - r}{w}\frac{Q}{L_r(k)}, & \text{если муравей } r \text{ прошёл через ребро } (v_i,jv_j), \\[1ex]
		0, & \text{иначе},
	\end{cases}
\end{equation}

Также как и в элитарном варианте дополнительно добавляется феромон по лучшему пути, давай в итоге формулу: 
\begin{equation}
	\tau_{ij}(k+1)=(1 - \rho)\tau_{ij}(k) + \sum_{r=1}^{w} \Delta \tau_{ij}^{(r)}(k) + \cdot \Delta\tau_{ij}^{best},
\end{equation}


\textbf{Улучшения относительно классического алгоритма}

\begin{itemize}
	\item ускоренная сходимость за счёт использования лучших путей;
	\item уменьшения шанса падения в локальный минимум, за счёт нескольких лучших путей, покрывающих граф шире;
\end{itemize}

\textbf{Ухудшения относительно классического алгоритма}

\begin{itemize}
	\item сложность подбора оптимального параметра $w$;
	\item дополнительные вычисления при поиске лучших путей в каждом поколении;
\end{itemize} 

\subsection{Гибридные алгоритмы}

Основной идее гибридных алгоритмов является комбинирование муравьиного алгоритма с другими эвристиками с целью закрытия их слабых мест друг другом. Механизмы их гибридизации могут различны, начиная с постобработки решений одного алгоритма другим, заканчивая ограничением выборов одного алгоритма другим.


\subsection*{Гибрид муравьиного алгоритма и табу-поиска (МАТП)}

Гибридизация муравьиного алгоритма с табу-поиском~\cite{aco-ts} (TS) строится на идее сочетания глобального поиска муравьиного алгоритма и локального поиска за счёт использования механизма краткосрочной памяти, характерного для табу поиска.

Табу-поиск работает следующим образом:

\begin{enumerate}
	\item начиная с некоторого заданного решения, генерируются соседние с ним решения из которых выбирается лучшее. В качестве соседних берутся полученные следующими операциями:
	\begin{itemize}
		\item 2-opt: удаление двух рёбер и пересоединение путей другим способом;
		\item 3-opt: удаление трёх рёбер и пересоединение;
		\item Swap (обмен): перестановка двух узлов в маршруте;
		\item Insertion: перемещение узлов в другую позицию.
	\end{itemize}
	\item в табу список заносится информация, обратная сделанному изменения. Таким образом запрещаются дальнейшие операции, которые приведут к уже исследованному решению;
	\item для текущего решения подбираются соседние с ним, из которых выбирается лучшее, но которое не противоречит списку табу.  Если существует хотя бы одно решение, не противоречащее списку табу, то выбирается лучшее из них, изменение заносится в список табу и алгоритм продолжается. Если не существует решений, непротиворечащих списку, то алгоритм заканчивается. На данном этапе возможно, что лучшее решение хуже текущего, тем не менее  оно все равно будет использовано, чтобы избегать локального минимума.
\end{enumerate}

Табу-поиск в данном гибридном варианте применяется ко всем путям, которые были построены муравьями, при этом феромоны обновляются уже по лучшим путям, полученным табу-поиском (если он стал лучше после табу-поиска).

\textbf{Улучшения относительно классического алгоритма}

\begin{itemize}
	\item улучшение качества найденных решений;
	\item ускорение сходимости, за счёт детерминированного выбор более коротких участков;
\end{itemize}

\textbf{Ухудшения относительно классического алгоритма}

\begin{itemize}
	\item увеличенная вычислительная сложность;
\end{itemize} 


\subsection*{Гибрид муравьиного алгоритма и алгоритма слизевика (МААС)}

В отличие от табу-поиска, где комбинация алгоритмов происходила за счёт локального улучшения решений построенных с помощью классического муравьиного алгоритма, данный вариант предлагает другую идею гибридизации алгоритмов.

Алгоритм слизевика (SMA)~\cite{smacfa} основан на биологическом феномене: простейший организм <<Physarum polycephalum>> способен находить кратчайшие пути в лабиринтах из точки его локализации в точку с пищей, распространяя себя как сеть трубочек, которые утолщаются, если они эффективнее остальных передают пищу.

Для решения задачи коммивояжёра слизевик имитируется как динамическая трубопроводная сеть, в которой каждому ребру графа сопоставлены несколько переменных, описывающих его физические характеристики.

Основными переменными модели являются:
\begin{itemize}
	\item $L_{ij}$ -- вес ребра $e_{ij}$ ;
	\item $D_{ij}(t)$ -- проводимость трубопровода на ребре $e_{ij}$ в момент времени $t$;
	\item $P_i$ -- <<давление>> в вершине $v_i$, формируемое в соответствии с законами сохранения потока;
	\item $Q_{ij}$-- поток вещества по ребру $e_{ij}$ , определяемый разностью давлений и проводимостью.
\end{itemize}

Имитация слизевика происходит через итеративное моделирование физических процессов в трубах согласно законам гидродинамики. При этом используется следующая последовательность действий:

\begin{enumerate}
	\item вычисляется давление во всех заданных узлах графа по закону Кирхгофа: $$\sum_i \frac{D_{ij}}{L_{ij}} (P_i - P_j) = S_j,$$
	где $S_j$ равна положительному потоку источника, если j -- начальная точка, отрицательному -- если j -- конечная точка и нулю в иных случаях.
	\item вычисляется поток по каждому ребру, согласно формуле: $$Q_{ij} = \frac{D_{ij}}{L_{ij}} (P_i - P_j).$$
	\item рассчитывается проводимость всех рёбер в следующий момент времени: $$D_{ij}(t+1) = D_{ij}(t) + \left(\frac{|Q_{ij}|}{1 + |Q_{ij}|} - D_{ij}(t)\right)\Delta t.$$
\end{enumerate}

Итерации продолжаются до тех пор, пока значение $D_{ij}$ не стабилизируется и изменения между шагами не станут меньше заданного порога. Лучшими узлами считаются те, которые имеют высокий показатели потока $Q_{ij}$ и проводимости $D_{ij}$. Лучший путь при этом строится жадным алгоритмом начиная с начальной точки, что приводит к тому, что плохо учитывается завершающий цикл переход.

При гибридном подходе, предложенном в~\cite{smacfa}, в начале алгоритма единожды просчитывается путь слизевика и в список $L_{AB}$ отбираются рёбра, которые имеют показатели потока $Q_{ij}$ и проводимости $D_{ij}$ выше заданного порога: $$Q_{ij} > Q_{min}, D_{ij} > D_{min}$$. 

Далее используется классический муравьиный алгоритм, но с модифицированным правилом выбора очередного перехода на пути муравья:

\begin{itemize}
	\item если среди возможных переходов есть тот, который есть в списке $L_{AB}$, то он выбирается безусловно;
	\item если такого перехода нет, то выбор проходит в соответствие с обычным вероятностным правилом муравьиного алгоритма.
\end{itemize}

\textbf{Улучшения относительно классического алгоритма}

\begin{itemize}
	\item улучшение качества найденных решений;
	\item ускорение сходимости, за счёт детерминированного выбора качественных частей;
\end{itemize}

\textbf{Ухудшения относительно классического алгоритма}

\begin{itemize}
	\item увеличение вычислительной сложности алгоритма, за счёт предрасчёта алгоритмом слизевика;
	\item увеличение вероятности падения в локальный минимум;
	\item существенное усложнение параметризации, из-за большого числа параметров.
\end{itemize} 

\subsection*{Гибрид муравьиного алгоритма и генетического (ГАМА)}

Генетический алгоритм (GA)~\cite{gaco} представляет собой метод оптимизации, основанный на принципах естественного отбора и наследования. Популяция возможных решений представляется в виде хромосом -- обычных фиксированных решений задачи, т.~е. в случае задачи коммивояжёра -- циклов в графе. Алгоритм ищет решения скрещивая лучшие решения в популяциях. Процесс поиска организован циклически и включается в себя:

\begin{enumerate}
	\item инициализация -- формирование начальной популяции хромосом;
	\item оценка приспособленности -- вычисление качества каждого решения по заданной целевой функции;
	\item отбор -- выбор части особей для размножения пропорционально их приспособленности.
	\item скрещивание -- обмен фрагментами хромосом между выбранными родителями. В случае задачи коммивояжёра решения обмениваются участками их циклов, из которых получается новое решение;
	\item мутация -- случайные изменения генов, которое нужно для исследования новых решений. В случае задачи коммивояжёра это может быть обмен позициями в цикле двух узлов.
\end{enumerate}

Процесс поиска, как и в муравьином алгоритме, проходит до определённого количества поколений или до найденного с заданной точностью решения.

Гибридный алгоритм ГАМА, предложенный в~\cite{gaco}, комбинирует муравьиный и генетический алгоритмы, при этом в отличие от предыдущих гибридных алгоритмов, в данном варианте генетический алгоритм выступает в роли ограничителя для каждого из муравьев.

Хромосома в данном случае не задаёт напрямую решение, а кодирует стратегию, по которой муравей должен выбирать следующий узел при построении пути. Хромосома представляет из себя непротиворечивые множество списков узлов, в которые муравей может отправиться на каждом шаге.

\begin{enumerate}
	\item Создаётся первая популяция муравьёв со случайными хромосомами;
	\item Муравьи строят маршруты с учётом ограничений из хромосомы:
	\begin{itemize}
		\item если для текущего шага у муравья всего один разрешённый узел, то он выбирает его для перехода;
		\item если таких узлов несколько, то выбор между ними делается по обычному вероятностному правилу муравьиного алгоритма.
	\end{itemize}
	\item Полученные решения оцениваются и из лучших происходит формирование следующего поколения хромосом за счёт скрещивания и мутации
	\item Происходит  обновление феромонов, как в классическом муравьином алгоритме.
\end{enumerate}

Таким образом, в ГАМА комбинируются два вида обучения:
\begin{itemize}
	\item усиление удачных рёбер муравьиным алгоритмом;
	\item улучшение стратегий выбора с сохранением разнообразия в маршрутах генетическим алгоритмом.
\end{itemize}

\textbf{Улучшения относительно классического алгоритма}

\begin{itemize}
	\item уменьшение вероятности падения в локальный минимум, за счёт разнообразия в хромосомах;
	\item улучшение сходимости, за счёт двойного эволюционного подхода.
\end{itemize}

\textbf{Ухудшения относительно классического алгоритма}

\begin{itemize}
	\item существенное увеличение вычислительной сложности, за счёт отбора лучших и скрещивания;
	\item усложнение параметризации, из-за большого числа параметров.
\end{itemize} 

\subsection{Алгоритмы с динамической параметризацией}

В классическом муравьином алгоритме все параметры -- $\alpha$, $\beta$, $p$, $m$ -- фиксированы и задаются до начала алгоритма, однако в процессе проведения расчётов важность факторов может меняться, из-за чего константные параметры могут быть неэффективными, в определённые моменты поиска решения. Алгоритмы данного раздела вводят способы, как изменять параметры по ходу расчёта, для увеличения эффективности или решения проблем муравьиного алгоритма.

\subsection*{Адаптивный муравьиный алгоритм с динамическими коэффициентами (АМАДК)}

В данном алгоритме, предложенном в~\cite{adla} вводится ряд изменений относительно классического муравьиного алгоритма:

\begin{itemize}
	\item вводятся формулы для динамического расчёта $\alpha$ и $\beta$ для каждой итерации алгоритма;
	\item вводятся стратегии по локальным улучшения полученных путей;
	\item изменяются способы обновления феромона на пройденных путях;
\end{itemize}

\textbf{Формулы для динамического расчёта $\alpha$ и $\beta$}

В данной модификации, в отличие от классического алгоритма, коэффициенты влияния феромона и эвристика рассчитываются динамически, для каждой итерации. При этом формулы их расчёта выбраны исходя из следующих тезисов: Вначале алгоритма концентрация низка на всех рёбрах, поэтому выбор прежде всего строится на эвристике, и, чтобы избежать выбора только относительно эвристики, важно поддерживать высокое значение $\alpha$ с относительно низким $\beta$. Однако с продолжением расчётов и накоплением феромонов, дабы избежать падения в локальный минимум, но продолжать искать лучшие пути нужно постепенно уменьшать значения $\alpha$ и увеличивать $\beta$. При этом для увеличения разнообразия, добавляется случайный множитель. Предложенные формулы расчёта имеют вид:
$$
\alpha(n_c)=\cos\left(\frac{r_1 n_c\pi}{2 n_{c\max}}\right)+A,
$$
$$
\beta(n_c)=\sin\left(\frac{r_2 n_c\pi}{2 n_{c\max}}\right)+B,
$$
где
\begin{itemize}
	\item $r_1, r_2 \in [0,1]$ -- случайные числа;
	\item $A=2, B=3$ -- экспериментально подобранные константы;  
	\item $n_c$  -- текущая итерация;
	\item $n_{c\max}$  -- максимальное число итераций.
\end{itemize}

\textbf{Адаптивный расчёт феромона}

Данный алгоритм проводит обновление феромона в 2 этапа: локальный и глобальный.

Локально, концентрация феромона обновляется после каждого прохода муравья. Формула для этого обновления:

$$
\tau_{ij}(t)= (1-\varepsilon)\tau_{ij}(t)+\frac{\varepsilon}{m\cdot L_{nn}},
$$

где
\begin{itemize}
	\item $\epsilon \in [0, 1]$ -- коэффициент испарения, для локального обновления;
	\item $L_{nn}$ -- длина пути, полученного жадным алгоритмом.
\end{itemize}

Глобальное обновление происходит также, как и в ранговом муравьином алгоритме, однако коэффициент испарения $p$ рассчитывается динамически для каждой итерации. По умолчанию он статичен и меняется в том случае, если лучшее решение не изменяется в течение нескольких итераций, что значит, что алгоритм попал в локальный минимум и коэффициент испарения нужно увеличить. Формула для изменения имеет следующий вид:

$$
\rho(n_c + 1) = 
\begin{cases}
	\rho_0, & n_c < \omega \cdot n_{c_{\max}} \\
	\frac{\rho(n_c)}{\gamma} , & n_c > \omega \cdot n_{c_{\max}} \text{ и } s > s_0,
\end{cases}
$$
где
\begin{itemize}
	\item $\rho_0$ -- начальное (максимальное) значение коэффициента испарения;
	\item $\omega \in [0, 1]$ -- множитель, определяющий после какого числа операций требуется проверка на локальный минимум;
	\item $s$ и $s_0$ -- количество итераций без улучшения результата и пороговое значение соответственно;
	\item $\gamma \in [0, 1]$ -- множитель коэффициента испарения.
\end{itemize}

Таким образом, данная модификация предлагает решения проблем стагнации и падения в локальный минимум за счёт динамических манипуляций с параметрами, отвечающими за феромон и влияние феромона на путь муравьёв;

\textbf{Улучшения относительно классического алгоритма}

\begin{itemize}
	\item улучшение качества результатов, за счёт использования более удачных параметров на разных этапах поиска;
	\item уменьшение вероятности падения в локальный минимум, за счёт эвристики в начале поиска и увеличение испарения при стагнации;
\end{itemize}

\textbf{Ухудшения относительно классического алгоритма}

\begin{itemize}
	\item чувствительность к параметрам в формулах выбора пути;
	\item эмпиричность формул, которые могут быть не эффективными для разных классов задач;
	\item существенное усложнение параметризации, из-за большого числа параметров.
\end{itemize} 

\subsection*{Муравьиный алгоритм с динамическим коэффициентом испарения и условием завершения (МАЭ)}

Данный алгоритм, описанный в~\cite{admut}, вводит понятие энтропии решений муравьёв, на основе которых в нём динамически меняются параметры и проверяется условие завершения алгоритма. Относительно классического муравьиного алгоритма, в данном варианте введены 3 изменения:

\begin{enumerate}
	\item кластеризация узлов;
	\item адаптивное испарение феромонов;
	\item условие завершения на основе энтропии.
\end{enumerate}

В классическом муравьином алгоритме при построении маршрута муравей выбирает следующий узел из всех еще не посещенных узлов, что приводит к вычислительной сложности $O(n^2)$ для построения маршрута одного муравья. В МАЭ для каждого узла создается набор кластеров, которые вместе содержат все остальные узлы. Кластеры упорядочены по вероятности того, что узел в кластере является частью оптимального маршрута. Наиболее вероятные узлы помещаются в первичные кластеры.

Процесс построения маршрута муравья с использованием кластеризации узлов включает два этапа:
\begin{enumerate}
	\item Выбор кластера: для последнего узла в текущем маршруте вычисляются вероятности выбора каждого кластера на основе средних значений эвристики и феромонов для узлов в кластере: $$
	p(C_j) = \frac{(\eta_{C_j})^\alpha \cdot (\tau_{C_j})^\beta}{\sum_{i=1}^{n_{clusters}} (\eta_{C_i})^\alpha \cdot (\tau_{C_i})^\beta}
	$$
	
	где $\eta_{C_j}$ - средняя эвристическая информация (обратные расстояния) между текущим узлом и узлами кластера $C_j$, $\tau_{C_j}$ - средняя интенсивность феромона.
	
	\item Выбор узла из кластера: если в выбранном кластере есть непосещенные узлы, то следующий узел выбирается из них классическим способом. Если в выбранном кластере нет непосещенных узлов, то используются следующие кластеры по порядку.
\end{enumerate}

\textbf{Адаптивное управление феромонами на основе энтропии}

Под энтропией в данной работе подразумевается степень разнообразия решений, полученных в рамках одного поколения. Для вычисления энтропии сначала определяется вероятность появления каждого ребра в решениях текущей популяции:

$$
p_{ij} = \frac{m_{ij}}{m},
$$
$$
H = -\sum_{i=1}^{n}\sum_{j=i+1}^{n} p_{ij} \cdot \log_2 p_{ij},
$$
где $m_{ij}$ -- количество муравьёв, прошедших через ребро $e_{ij}$.


Высокие значения энтропии соответствуют ситуации, когда муравьи используют разнообразные наборы рёбер, что характерно для начала алгоритма. Низкие значения энтропии указывают на концентрацию решений вокруг небольшого набора рёбер, что происходит при сходимости к локальному оптимуму. Относительно энтропии рассчитывается коэффициент испарения для каждого поколения:
$$
\rho = \rho_{min} + (\rho_{max} - \rho_{min}) \cdot \frac{H - H_{min}}{H_{max} - H_{min}}
$$

\begin{itemize}
	\item при высокой энтропии $\rho \rightarrow \rho_{max}$, что ускоряет испарение и стимулирует исследование новых областей пространства поиска;
	\item при низкой энтропии $\rho \rightarrow \rho_{min}$, что замедляет испарение и позволяет углубить поиск более оптимального локального решения.
\end{itemize}


\textbf{Условие завершения на основе энтропии}

Традиционные критерии остановки (фиксированное число итераций или отсутствие улучшений) часто оказываются неоптимальными: либо преждевременно останавливают поиск, либо приводят к избыточным вычислениям.

Новое условие завершения использует энтропию как индикатор исчерпания потенциала улучшений:
$$
H \cdot (1 + \omega) \leqslant H_{min} 
$$
где $\omega$ - малый положительный параметр, определяющий допустимую близость к минимальной энтропии.

\textbf{Улучшения относительно классического алгоритма}

\begin{itemize}
	\item оптимизация числа итераций, за счёт динамического условия завершения;
	\item уменьшение вероятности падения в локальный минимум, так как коэффициент испарения увеличивается при низкой энтропии, что способствует разнообразию решения;
\end{itemize}

\textbf{Ухудшения относительно классического алгоритма}

\begin{itemize}
	\item увеличение вычислительной сложности, за счёт просчёта энтропии на каждой итерации;
	\item замедление сходимости, за счёт динамического испарения.
\end{itemize}

\subsection*{Муравьиный алгоритм с динамическими гетерогенными параметрами (ГДМА)}

В данном алгоритме, предложенном в~\cite{hetero}, реализуется идея гетерогенных, т.~е. разнородных параметров муравьёв. Основная идея заключается в том, что каждый муравей в колонии обладает индивидуальными поведенческими характеристиками, которые эволюционируют в процессе поиска решения.

В отличие от классического подхода, где все муравьи используют одинаковые значения параметров $\alpha$ и $\beta$, в ГДМА каждый муравей $k$ имеет собственные значения $\alpha_k$ и $\beta_k$, определяющие его предпочтения при выборе пути:

$$
P^{k}_{ij} = \frac{[\tau_{ij}]^{\alpha_k} [\eta_{ij}]^{\beta_k}}{\sum_{l \in N^k} [\tau_{il}]^{\alpha_k} [\eta_{il}]^{\beta_k}}
$$

где $N^k$ - множество допустимых переходов для муравья $k$.

Параметры $\alpha$ и $\beta$ у муравьёв адаптируются с использованием упрощённого эволюционного алгоритма, при этом процесс можно разделить на несколько итераций:
\begin{enumerate}
	\item \textbf{инициализация}: по умолчанию все муравьи инициализируются некоторыми параметрами, при этом в самой~\cite{hetero} предложен вариант с жадными муравьями, где $\alpha_k = 0$ и $\beta_k = 10$;
	\item \textbf{оценка эффективности}: каждые $w$ итераций вычисляется средняя приспособленность каждого муравья на основе длины построенных им маршрутов. 
	\item \textbf{селекция}: выбирается муравей с наилучшей средней приспособленностью и муравей с наихудшей средней приспособленностью.
	\item \textbf{мутация}: параметры лучшего муравья подвергаются гауссовской мутации для создания нового потомка:
	$$
	\alpha_{child} = \alpha_{best} + \mathcal{N}(0, \sigma), \quad \beta_{child} = \beta_{best} + \mathcal{N}(0, \sigma)
	$$
	где $\sigma = 0.05$ - стандартное отклонение.
	\item \textbf{замена}: созданный потомок заменяет наихудшего муравья в популяции.
\end{enumerate}

\textbf{Улучшения относительно классического алгоритма}

\begin{itemize}
	\item уменьшения вероятности падения в локальный минимум, за счёт разнообразия параметров у муравьёв;
	\item улучшение качества решения, за счёт подстраивания параметров, под разные этапы поиска;
	\item упрощение параметризации, меньшая чувствительность к начальным значениям, за счёт подбора параметров динамически.
\end{itemize}

\textbf{Ухудшения относительно классического алгоритма}

\begin{itemize}
	\item увеличенная чувствительность к количеству итераций, так как подбор оптимальных параметров происходит в зависимости от входного графа;
	\item замедленная сходимость, за счёт подбора параметров.
\end{itemize}

\section{Сравнение рассмотренных методов}

Как уже было указано ранее, модификации муравьиного алгоритма и сам алгоритм, состоят из 3-х структурных частей:
\begin{itemize}
    \item стратегия выбора пути;
    \item правила обновления феромона;
    \item критерии остановки.
\end{itemize}

Каждая модификация изменяет один или несколько механизмов, что и описывает её ключевые отличия от муравьиного алгоритма. Также в модификация меняется количество параметров, что влияет на сложность предварительной настройки. Эти пункты позволяют кратко описать основные новшества, введённые модификацией.

Любая модификация алгоритма стремится решить один или несколько недостатков классического, при этом зачастую жертвуя чем-то другим, поэтому важно указывать, что именно улучшила модификация, и чем ради этого пришлось пожертвовать.

В таблице~\ref{tbl:criteria} представлены сформулированные критерии сравнения рассмотренных модификаций муравьиного алгоритма.

\begin{longtable}{|c|p{0.2\textwidth}|p{0.65\textwidth}|}
	\caption{Критерии сравнения модификаций муравьиного алгоритма}\label{tbl:criteria} \\
	\hline
	\textbf{№} & \textbf{Критерий} & \textbf{Описание} \\
	\hline
	\endfirsthead
	\caption[]{Критерии сравнения модификаций муравьиного алгоритма (продолжение)} \\
	\hline
	\textbf{№} & \textbf{Критерий} & \textbf{Описание} \\
	\hline
	\endhead
	\hline
	\endfoot
	\endlastfoot
	1 & Число гиперпараметров & Количество настраиваемых параметров, которые необходимо задать до начала работы алгоритма. \\
	\hline
	2 & Стратегия построения цикла & Принцип, по которому отдельный муравьёв строит полный маршрут (тур). \\
	\hline
	3 & Стратегия обновления феромона & Правила изменения уровня феромона на рёбрах графа (испарение и добавление). \\
	\hline
	4 & Критерий остановки & Условие, при выполнении которого работа алгоритма завершается. \\
	\hline
	5 & Улучшения & Улучшения, привнесённые модификацией по сравнению с классическим муравьиным алгоритмом. \\
	\hline
	6 & Ухудшения & Недостатки или компромиссы модификации по сравнению с классическим муравьиным алгоритмом. \\
	\hline
\end{longtable}

В таблицах~\ref{tbl:comparison-basic}-~\ref{tbl:comparison-adaptive} представлено сравнение модификаций по сформулированным критериям.


\begin{longtable}{|p{0.02\textwidth}|p{0.27\textwidth}|p{0.27\textwidth}|p{0.27\textwidth}|}
	\caption{Сравнение расширений классического муравьиного алгоритма}\label{tbl:comparison-basic} \\
	\hline
	\textbf{№} & \textbf{ЭМА} & \textbf{МММА} & \textbf{РМА} \\
	\hline
	\endfirsthead
	\caption[]{Сравнение расширений классического муравьиного алгоритма (продолжение)} \\
	\hline
	\textbf{№} & \textbf{ЭМА} & \textbf{МММА} & \textbf{РМА} \\
	\hline
	\endhead
	\hline
	\endfoot
	\endlastfoot
	1 & $6$ ($\alpha, \beta, \rho, m, Q, e$) & $7$ (+ $\tau_{min}, \tau_{max}$) & $6$ ($\alpha, \beta, \rho, m, Q, w$) \\
	\hline
	2 & Стандартная & Стандартная & Стандартная \\
	\hline
	3 & Все муравьи + $e$ по лучшему пути & Только лучший муравей. $\tau_{ij} \in [\tau_{min}, \tau_{max}]$ & $w$ лучших (взвеш. по рангу) + лучший путь \\
	\hline
	4 & Стандартные & Стандартные & Стандартные \\
	\hline
	5 & \begin{itemize}[leftmargin=*, nosep]
		\item Ускоренная сходимость
	\end{itemize} & \begin{itemize}[leftmargin=*, nosep]
		\item Ускоренная сходимость
		\item Снижен риск лок. минимума
	\end{itemize} & \begin{itemize}[leftmargin=*, nosep]
		\item Ускоренная сходимость
		\item Снижен риск лок. минимума
	\end{itemize} \\
	\hline
	6 & \begin{itemize}[leftmargin=*, nosep]
		\item Риск стагнации/лок. мин.
	\end{itemize} & \begin{itemize}[leftmargin=*, nosep]
		\item Сложная параметризация ($\tau_{min}, \tau_{max}$)
	\end{itemize} & \begin{itemize}[leftmargin=*, nosep]
		\item Сложный подбор $w$
		\item Доп. вычисления (ранжирование)
	\end{itemize} \\
	\hline
\end{longtable}

\begin{longtable}{|p{0.02\textwidth}|p{0.27\textwidth}|p{0.27\textwidth}|p{0.27\textwidth}|}
	\caption{Сравнение гибридных модификаций муравьиного алгоритма}\label{tbl:comparison-hybrid} \\
	\hline
	\textbf{№} & \textbf{МАТП} & \textbf{МААС} & \textbf{ГАМА} \\
	\hline
	\endfirsthead
	\caption[]{Сравнение гибридных модификаций муравьиного алгоритма (продолжение)} \\
	\hline
	\textbf{№} & \textbf{МАТП} & \textbf{МААС} & \textbf{ГАМА} \\
	\hline
	\endhead
	\hline
	\endfoot
	\endlastfoot
	1 & $5$ ($\alpha, \beta, \rho, m, Q$) & $8$ ($\alpha, \beta, \rho, m, Q,$ $Q_{min}, D_{min}, S$) & $7$ ($\alpha, \beta, \rho, m, Q$, параметры ГА) \\
	\hline
	2 & Стандартный + лок. оптимизация & Безусловный выбор рёбер из пути слизевика, иначе стандартный & Стандартный, но из множества, разрешённого хромосомой \\
	\hline
	3 & По улучшенным путям & Стандартный & Стандартный \\
	\hline
	4 & Стандартные & Стандартные & Стандартные \\
	\hline
	5 & \begin{itemize}[leftmargin=*, nosep]
		\item Качество решений
		\item Ускоренная сходимость
	\end{itemize} & \begin{itemize}[leftmargin=*, nosep]
		\item Качество решений
		\item Ускоренная сходимость
	\end{itemize} & \begin{itemize}[leftmargin=*, nosep]
		\item Устойчивость к лок. мин.
		\item Улучшенная сходимость
	\end{itemize} \\
	\hline
	6 & \begin{itemize}[leftmargin=*, nosep]
		\item Высокая выч. сложность
	\end{itemize} & \begin{itemize}[leftmargin=*, nosep]
		\item Высокая выч. сложность
		\item Риск лок. минимума
		\item Сложная параметризация
	\end{itemize} & \begin{itemize}[leftmargin=*, nosep]
		\item Высокая выч. сложность
		\item Сложная параметризация
	\end{itemize} \\
	\hline
\end{longtable}



\begin{longtable}{|p{0.02\textwidth}|p{0.27\textwidth}|p{0.27\textwidth}|p{0.27\textwidth}|}
	\caption{Сравнение динамических модификаций муравьиного алгоритма}\label{tbl:comparison-adaptive} \\
	\hline
	\textbf{№} & \textbf{АМАДК} & \textbf{МАЭ} & \textbf{ГДМА} \\
	\hline
	\endfirsthead
	\caption[]{Сравнение динамических модификаций муравьиного алгоритма (продолжение)} \\
	\hline
	\textbf{№} & \textbf{АМАДК} & \textbf{МАЭ} & \textbf{ГДМА} \\
	\hline
	\endhead
	\hline
	\endfoot
	\endlastfoot
	1 & $10$ ($\alpha, \beta, \rho, m, Q,$ $A, B, \omega, \gamma, s_0$) & $7$ ($\alpha, \beta, \rho, m, Q, \rho_{min}, \rho_{max}$, $H_{min}$, $\omega$) & $7$ ($\alpha, \beta, \rho, m, Q, \sigma$, $w_{adapt}$) \\
	\hline
	2 & Стандартный + лок. улучшения & Кластеризация узлов $\rightarrow$ выбор внутри кластера & Стандартный, но с индивид. параметрами \\
	\hline
	3 & Ранговое + испарение растёт при стагнации & Испарение повышается при низкой энтропии решений & Стандартный \\
	\hline
	4 & Стандартный & Энтропия меньше заданного порога & Стандартный \\
	\hline
	5 & \begin{itemize}[leftmargin=*, nosep]
		\item Качество решений
		\item Устранение стагнации
	\end{itemize} & \begin{itemize}[leftmargin=*, nosep]
		\item Оптим. число итераций
		\item Устранение стагнации
	\end{itemize} & \begin{itemize}[leftmargin=*, nosep]
		\item Устранение стагнации
		\item Качество решений
		\item Проще параметризация
	\end{itemize} \\
	\hline
	6 & \begin{itemize}[leftmargin=*, nosep]
		\item Эмпиричность формул
		\item Сложная параметризация
	\end{itemize} & \begin{itemize}[leftmargin=*, nosep]
		\item Замедленная сходимость
		\item Выч. сложность ($H$, кластеры)
	\end{itemize} & \begin{itemize}[leftmargin=*, nosep]
		\item Замедленная сходимость
		\item Чувств. к числу итераций
	\end{itemize} \\
	\hline
\end{longtable}

Проведённый сравнительный анализ модификаций муравьиного алгоритма для задачи коммивояжёра показывает, что каждая группа модификаций предлагает свой компромисс между качеством решения, скоростью сходимости, вычислительной сложностью и трудоёмкостью настройки. Базовые расширения (ЭМА, МММА, РМА) фокусируются на совершенствовании механизма феромонов, что улучшает сходимость, но не решает фундаментальных проблем стагнации и чувствительности к параметрам. Гибридные подходы (ACO-TS, МААС, ГАМА), интегрирующие внешние эвристики, демонстрируют значительный потенциал в повышении качества решений, однако зачастую ценой существенного усложнения алгоритма и роста вычислительных затрат.

Наиболее перспективными представляются алгоритмы с динамической параметризацией, так как они позволяют уменьшить зависимость алгоритма от входных данных. Особенно выделяется эволюционный подход к реализации подбора параметров, так как он позволяет динамически менять параметры в зависимости от конфигурации графа и процесса поиска решения. При этом наличие индивидуальных параметров позволяет поддерживать разнообразие решения, направляя его в наиболее эффективное русло. Рассмотренные выше модификации изменяют параметры с помощью эмпирических формул или случайного перебора, не учитывая динамическую среду поиска, что ограничивает их возможность адаптивно подстраиваться под нее.Динамический подбор параметров можно оптимизировать путём использования алгоритмов, адаптирующих поведение муравьёв в зависимости от качества их решений.

Таким образом, объединение муравьиного алгоритма с более сложными эволюционным алгоритмом, таким как генетический, имеет высокие шансы оказаться более эффективным других подходов в неспецифичных случаях. Одним из вариантов дальнейшего развития муравьиного алгоритма предлагается <<Метод решения задачи коммивояжёра муравьиным алгоритм с гетерогенными динамическими параметрами на основе генетического алгоритма>>, формализация которого в нотации IDEF-0 представлена на рисунке~\ref{fig:diplom}.

\begin{figure}[H]
	\centering
	\includegraphics[width=0.9\textwidth]{method/01\_A0}
	\caption{Формализация предлагаемого метода в виде IDEF0-диаграммы}
	\label{fig:diplom}
\end{figure}

\section*{Вывод}


В результате аналитической части была формализована задача коммивояжёра в терминах теории графов и нотации IDEF0, описан классический муравьиный алгоритм и выявлены его недостатки. Были рассмотрены и классифицированы существующие модификации муравьиного алгоритма, проведён их сравнительный анализ по ключевым критериям. На основе анализа обоснована перспективность использования эволюционного подхода для динамического подбора гетерогенных параметров, что позволило сформулировать предлагаемый метод решения задачи коммивояжёра муравьиным алгоритмом с гетерогенными динамическими параметрами на основе генетического алгоритма.
